%% Unary normal modal Lindenbaum algebras and general-frame completeness.

\subsection{Normal modal logics, Lindenbaum algebras, and general frames}
\label{sec:p10-modal-algebra-completeness}

\begin{hbdefinition}{def:p10-normal-modal-algebra-logic}{The fixed modal language and its algebras}
Fix the propositional-variable set \(\mathsf{Prop}\) of Part~VIII.
Formulas are built using \(\bot,\neg,\wedge,\Box\), with the other
Boolean connectives abbreviated as usual and
\(\Diamond\varphi=\neg\Box\neg\varphi\).
A \emph{normal modal logic} \(L\) is a set of formulas containing
all classical propositional tautology instances and
\[
 \Box(\varphi\to\psi)\to(\Box\varphi\to\Box\psi),
\]
and closed under modus ponens, uniform substitution, and
theorem necessitation: \(\varphi\in L\) implies \(\Box\varphi\in L\).
Write \(\vdash_L\varphi\) for \(\varphi\in L\).
The logic is consistent when \(\bot\notin L\).
For premises \(\Gamma\), a derivation uses members of \(\Gamma\),
members of \(L\), and modus ponens.  Necessitation applies only
to theorems, whose boxes already belong to \(L\); uniform
substitution is likewise a closure rule for \(L\), not a rule
for freely substituting in premises.
This extends the local convention of
Definition~\ref{def:p08-modal-derivability} to an arbitrary normal \(L\).

A \emph{normal modal algebra} is a Boolean algebra \(A\), possibly
trivial, with a unary operation \(\Box_A\) satisfying
\[
 \Box_A1=1,\qquad
 \Box_A(a\wedge b)=\Box_Aa\wedge\Box_Ab.
\]
Its dual \(\Diamond_Aa=\neg\Box_A\neg a\) preserves \(0\) and
binary joins, so this is the unary case of a Boolean algebra
with operators.
A valuation \(v:\mathsf{Prop}\to A\) extends recursively to
\(\llbracket\varphi\rrbracket_v^A\).
Write \(A\models\varphi\) when that value is \(1_A\) for every
valuation, and let
\[
 \operatorname{Alg}(L)
   =\{A:A\text{ is a normal modal algebra and }
                   A\models\lambda\text{ for every }\lambda\in L\}.
\]
Thus membership in \(\operatorname{Alg}(L)\) requires validity
under all algebra valuations.
\end{hbdefinition}

\begin{hbdefinition}{def:p10-modal-lindenbaum}{The normal modal Lindenbaum algebra}
Put
\[
 \varphi\equiv_L\psi
       \quad\Longleftrightarrow\quad
       \vdash_L\varphi\leftrightarrow\psi.
\]
The following theorem verifies that this is a congruence.
Let \(F_L\) be the quotient of the formula algebra by this relation,
with classes \([\varphi]_L\) and induced operations
\[
\begin{gathered}
 0=[\bot]_L,\quad 1=[\top]_L,\quad
 \neg[\varphi]_L=[\neg\varphi]_L,\\
 [\varphi]_L\wedge[\psi]_L=[\varphi\wedge\psi]_L,\qquad
 \Box[\varphi]_L=[\Box\varphi]_L.
\end{gathered}
\]
This is the \emph{normal modal Lindenbaum algebra} of \(L\).
\index{Lindenbaum--Tarski algebra!normal modal logic}
\end{hbdefinition}

\begin{hbtheorem}{thm:p10-modal-algebra-completeness}{proved=annex:p10.modal-algebra.1}{Algebraic completeness for every normal modal logic}
The relation \(\equiv_L\) is a congruence, \(F_L\) is a normal
modal algebra, and \(F_L\in\operatorname{Alg}(L)\).
For every formula \(\varphi\),
\[
 \vdash_L\varphi
 \quad\Longleftrightarrow\quad
 [\varphi]_L=1
 \quad\Longleftrightarrow\quad
 F_L\models\varphi
 \quad\Longleftrightarrow\quad
 \operatorname{Alg}(L)\models\varphi.
\]
If \(\varphi\notin L\), the valuation \(p\mapsto[p]_L\)
refutes \(\varphi\) in \(F_L\).
If \(L\) is inconsistent, it contains every formula and
\(\operatorname{Alg}(L)\) consists exactly of the trivial modal
algebras.
\end{hbtheorem}

The quotient and algebraic completeness are the unary normal case of
\cite[Definitions~3.14 and~3.16, Proposition~3.15, and
Theorems~3.13 and~3.17, pp.~10--11]{Venema2005Preprint}.
For an arbitrary valuation in the quotient, the annex chooses
representatives only for the finitely many variables in the formula
being evaluated.

\begin{hbdefinition}{def:p10-modal-canonical-general-frame}{The canonical general frame and its underlying frame}
Suppose \(L\) is consistent.
Apply the unary case of Theorem~\ref{thm:p10-jonsson-tarski} to
\(F_L\), using the Boolean prime ideal theorem.
Let \(X_L\) be the ultrafilters of its Boolean reduct, and put
\[
 \widehat a=\{U\in X_L:a\in U\},\qquad
 U R_L V
   \quad\Longleftrightarrow\quad
       \forall a\in V\ \Diamond a\in U.
\]
Equivalently,
\[
 U R_L V
   \quad\Longleftrightarrow\quad
       \{a:\Box a\in U\}\subseteq V.
\]
Define
\[
 \mathfrak G_L=(X_L,R_L,\mathcal A_L),\qquad
 \mathcal A_L=\{\widehat a:a\in F_L\},\qquad
 \mathfrak F_L=(X_L,R_L).
\]
The admissible sets \(\mathcal A_L\) are exactly the Stone clopens
and are closed under Boolean and modal operations, as in
Definition~\ref{def:p10-bao-general-frame}.
The \emph{canonical general frame} \(\mathfrak G_L\) restricts
variable valuations to \(\mathcal A_L\).
Its \emph{underlying canonical frame} \(\mathfrak F_L\) instead
permits every subset of \(X_L\) as a variable value.
These are the unary constructions of
\cite[Example~5.13 and Definition~5.24, pp.~20, 22]{Venema2005Preprint}.
\end{hbdefinition}

\begin{hbtheorem}{thm:p10-modal-general-frame-completeness}{proved=annex:p10.modal-algebra.2}{General-frame completeness and the canonical truth lemma}
Assume the Boolean prime ideal theorem.
For consistent normal \(L\), the frame \(\mathfrak G_L\) validates
every member of \(L\) under every admissible valuation.
For the canonical valuation \(v_0(p)=\widehat{[p]_L}\),
\[
 \mathfrak G_L,v_0,U\models\varphi
           \quad\Longleftrightarrow\quad[\varphi]_L\in U.
\]
Consequently \(\mathfrak G_L\models\varphi\) exactly when
\(\varphi\in L\).
The map
\[
 U\longmapsto\{\varphi:[\varphi]_L\in U\}
\]
identifies \(\mathfrak F_L\) with the maximal-consistent-set
canonical frame of Method~\ref{def:p08-modal-canonical-model}.
\end{hbtheorem}

\begin{hbtheorem}{thm:p10-modal-general-frame-strong}{proved=annex:p10.modal-algebra.3}{Strong completeness for local general-frame consequence}
Assume the Boolean prime ideal theorem.
Let \(\operatorname{Gen}(L)\) be all nonempty general frames
validating every member of \(L\) under every admissible valuation.
For every set \(\Gamma\) of formulas,
\[
 \Gamma\vdash_L\varphi
       \quad\Longleftrightarrow\quad
 \Gamma\models^\ell_{\operatorname{Gen}(L)}\varphi.
\]
Here derivations have the theorem-only necessitation convention of
Definition~\ref{def:p08-modal-derivability}, and the right side
requires all premises at the same point.
For consistent \(L\), every failure of derivability already has
a counterexample in \(\mathfrak G_L\) under \(v_0\).
For inconsistent \(L\), \(\operatorname{Gen}(L)\) is empty and
both sides hold for all \(\Gamma,\varphi\).
\end{hbtheorem}

\begin{hbdefinition}{def:p10-modal-canonicity}{Canonicity and unrestricted valuations}
For a consistent normal logic in the fixed language, the
\emph{canonical-frame validity condition} is
\[
 \mathfrak F_L\models\lambda\quad\text{for every }\lambda\in L,
\]
where valuations range over all subsets of \(X_L\).
This is an additional condition beyond validity on
\(\mathfrak G_L\).

When varying the set of variables to a set \(Y\), interpret \(L\)
schematically: let \(L(Y)\) be the least normal logic over \(Y\)
containing all substitution instances of members of \(L\).
In the all-variable-set convention, \(L\) is \emph{canonical}
when the canonical frame of \(L(Y)\) validates \(L(Y)\) for
every \(Y\).
This is the convention in
\cite[Definition~6.10, p.~31]{Venema2005Preprint}.
The next theorem needs only the fixed-language condition.
\index{canonicity!normal modal logic}
\end{hbdefinition}

\begin{hbtheorem}{thm:p10-modal-canonicity-completeness}{proved=annex:p10.modal-algebra.4}{Completeness from canonical-frame validity}
Assume the Boolean prime ideal theorem.
If consistent normal \(L\) satisfies the canonical-frame validity
condition, then \(L\) is strongly complete for local consequence
over the class \(\operatorname{Fr}(L)\) of nonempty ordinary
Kripke frames validating \(L\).

A sufficient algebraic condition is that
\(\operatorname{Alg}(L)\) be closed under
\[
 A\longmapsto A^\sigma
       =\text{the full complex algebra of the ultrafilter frame of }A.
\]
This closure property is called canonicity of the algebra class.
\end{hbtheorem}

The representation embedding \(A\hookrightarrow A^\sigma\)
preserves evaluations whose variable values lie in its image.
Equations true in a subalgebra need not hold under all valuations
in the larger algebra.  Thus the embedding alone supplies
truth preservation on admissible values.  Together with the
Lindenbaum construction this gives general-frame completeness;
the additional validity or
closure condition above supplies the ordinary-frame conclusion.
The algebraic closure formulation is
\cite[Definition~6.9 and Theorem~6.12, p.~31]{Venema2005Preprint}.
The equivalence with logical canonicity in that source's
Theorem~6.11 uses its all-variable-set convention.
The Kripke-incomplete normal logics of
Theorem~\ref{thm:p08-kripke-incomplete-logics} illustrate why the
ordinary-frame conclusion needs its own hypotheses.
